\newpage
\titledquestion{Balancing Act}

\textbf{Balancing Act}

A sequence of parentheses is \textit{balanced} if every open parenthesis is
matched by a later close parenthesis, and no close parenthesis appears before
its matching open. We represent such a sequence as a list:

\begin{codeblock}
  datatype paren = L | R

  type parens = paren list
\end{codeblock}

where \code{L} stands for \code{(} and \code{R} stands for \code{)}.

So, for instance:
\begin{itemize}
  \item \code{[L, R]} is balanced, representing \code{()}
  \item \code{[L, L, R, R]} is balanced, representing \code{(())}
  \item \code{[L, R, L, R]} is balanced, representing \code{()()}
  \item \code{[]} is balanced, representing the empty sequence
  \item \code{[R, L]} is \textbf{not} balanced, representing \code{)(}
  \item \code{[L, L, R]} is \textbf{not} balanced, representing \code{((}\code{)}
\end{itemize}

\begin{parts}
  \part[3]\hfill

    A student proposes the following approach: ``count the number of \code{L}s
    and the number of \code{R}s, and return \code{true} if they are equal.''

    Give a value of type \code{parens} which shows this approach is wrong, and
    state in one sentence the property that this approach fails to capture.

    \begin{solutionorbox}[8em]
      \code{[R, L]} has one \code{L} and one \code{R}, but is not balanced.

      Counting ignores \textbf{order}: at no point may the number of \code{R}s
      seen so far exceed the number of \code{L}s seen so far. Equal totals is
      necessary, but not sufficient.
    \end{solutionorbox}

  \part[3]\hfill

    A second student, having read your answer above, proposes: ``keep a running
    counter, starting at 0. Add 1 for each \code{L}, subtract 1 for each
    \code{R}. Return \code{true} if the final counter is 0.''

    Give a value of type \code{parens} which shows that \textbf{this} approach
    is also wrong.

    \begin{solutionorbox}[8em]
      \code{[L, R, R, L]}. The counter goes \code{1, 0, ~1, 0}, ending at 0,
      but the sequence is not balanced --- it dips below zero at the third
      element.
    \end{solutionorbox}

  \newpage

  \part[10]\hfill

    Now, implement \code{balanced}:

    \spec
      {balanced}
      {\code{parens -> bool}}
      {\code{true}}
      {\code{balanced ps} $\eeq$ \code{true} if and only if \code{ps} is
      balanced, and \code{false} otherwise}

    \begin{constraint}
      Your solution must traverse the list \textbf{exactly once}, and that
      traversal must be a single call to \code{foldl}. You may not use
      explicit recursion, \code{rev}, or any auxiliary list.
    \end{constraint}

    \textbf{Hint:} Your two counterexamples above tell you two different things
    that can go wrong. What does your accumulator need to remember in order to
    detect \textit{both} of them?

    \begin{solutionorbox}[26em]
      The accumulator must track the running depth \textit{and} whether the
      sequence has already gone irrecoverably wrong. \code{int option} does
      this: \code{SOME d} means ``still valid, at depth \code{d}'', and
      \code{NONE} means ``already broken, and nothing can fix it''.

      \begin{codeblock}
        fun balanced ps =
          case
            foldl
              (fn (_, NONE) => NONE
                | (L, SOME d) => SOME (d + 1)
                | (R, SOME d) =>
                    if d = 0 then NONE else SOME (d - 1))
              (SOME 0)
              ps
          of
            SOME 0 => true
          | _ => false
      \end{codeblock}

      Note that the post-processing step is essential: ending in
      \code{SOME d} for \code{d > 0} means unmatched open parens remain, as in
      \code{[L, L, R]}.
    \end{solutionorbox}

  \newpage

  \part[4]\hfill

    Suppose we now want to handle two \textit{kinds} of bracket, which must be
    matched with their own kind:

    \begin{codeblock}
      datatype kind = Round | Square
      datatype bracket = Open of kind | Close of kind
    \end{codeblock}

    So \code{[Open Round, Open Square, Close Square, Close Round]} (that is,
    \code{([])}) is balanced, but
    \code{[Open Round, Open Square, Close Round, Close Square]} (that is,
    \code{([)]}) is not.

    You do \textbf{not} need to write any code. State what the accumulator type
    must become, and explain in one or two sentences why an \code{int} depth is
    no longer sufficient.

    \begin{solutionorbox}[14em]
      The accumulator must become something like \code{kind list option} ---
      that is, a \textbf{stack} of the currently-open bracket kinds (still
      wrapped in an \code{option}, or with a dedicated failure representation).

      A depth counter records only \textit{how many} brackets are open, not
      \textit{which} ones. On seeing \code{Close Round} we must check that the
      most recently opened bracket was a \code{Round}, which requires
      remembering the kinds in order.

      Observe that the previous part's \code{int} accumulator was really this
      stack all along --- with only one kind of bracket, the only information
      a stack carries is its length.
    \end{solutionorbox}
\end{parts}
